\section{Communication \& Manuel}

\subsection{Etape du programme \& Mode d'emploi utilisateur}

\noindent Notre application fonctionne en 4 étapes :


1- Sélection de l'image à traiter.

2- Extraction des étoiles présentes sur l'image

3- Reconnaissance de la constellation avec le set d'étoiles

4- Affichage de la constellation et de son histoire

\noindent Naturellement, il est nécessaire de faire communiquer tous les programmes responsables, surtout que notre projet est multilangage (Java et Python).

\vspace{0.2cm}

\noindent L'utilisateur lui, n'a que trois étape à suivre :

1- Cliquer sur Recognition pour sélectionner l'image de son choix

2 - Une fois le chargement fini (repéré par la disparition de la boussole centrale), il peut cliquer sur le bouton Constellation qui apparaît aussitôt

3- La Constellation apparaît et un texte avec, il peut fermer le logiciel ou recliquer sur Constellation pour revenir à l'écran principal

\subsection{Communications interprocessus}

\subsubsection{Interface | Détecteur d'étoiles}

\noindent \textbf{Java -> Python : Fichier txt}

\vspace{0.2cm}

\noindent La première étape de fonctionnement de l'application est la sélection de la photo du ciel étoilé.
\noindent L'interface doit donc communiquer l'emplacement de l'image pour qu'elle soit traitée par le détecteur d'étoiles.
\noindent Pour cela, on écrit le path de l'image choisi (via l'interface) dans un fichier txt nommé output.txt.


\subsubsection{Détecteur d'étoiles | Reconnaissance}

\noindent \textbf{Python -> Java : Fichier CSV}

\vspace{0.2cm}

\noindent La seconde étape est le traitement de l'image et l'extraction des étoiles. Pour fournir leur localisation à la partie reconnaissance, nous avons décider, après des tests infructueux avec le format json, d'utiliser le format \textbf{csv}.

\subsubsection{Reconnaissance | Interface}

\noindent \textbf{Java -> Java : Fichier txt}

\vspace{0.2cm}

\noindent Tout comme le passage de la première à la seconde étape, on peut facilement de communiquer le nom de la constellation reconnue. L'idée est encore une fois d'écrire le nom dans un \textbf{.txt}.